\addcontentsline{toc}{section}{CAPÍTULO XIII}
\chapter{CONCLUSIONES}

\section{Validación Preliminar del MVP}

La hipótesis de valor del MVP, que un sistema de dos fases (reglas contractuales + LLM) puede clasificar las liquidaciones hospitalarias con suficiente precisión para reducir el tiempo de auditoría en un 40 \%, recibió evidencia favorable pero no fue validada en producción. La arquitectura es viable (0 errores de ejecución, 0 errores de \textit{parsing} en 101 liquidaciones),

La tasa de recomendación de \textit{fast-track} alcanzó el 65.3 \%, por encima del rango meta de 50– 60 \%. Una tasa más alta de la esperada podría indicar que el umbral de decisión es demasiado permisivo y que se están dejando pasar casos que merecerían revisión. Dicho esto, hay dos razones para no alarmarse todavía. Por un lado, el resultado está lejos del techo de alerta del 75 \% que definimos en §9.3 como señal de calibración insuficiente. Por otro, ninguno de los 17 casos del \textit{Golden Set} que debían escalarse fue clasificado erróneamente como fast-track, es decir, la tasa de falsos negativos observada es cero. El problema es que 17 casos son pocos para hacer una afirmación fuerte al respecto, y además esos mismos casos se usaron para calibrar los umbrales, con lo cual hay un sesgo inherente. En la práctica, la pregunta de si el 65.3 \% es seguro o riesgoso solo puede responderla un piloto con auditores reales procesando lotes de producción. Lo que sí podemos afirmar es que el exceso sobre la meta se explica principalmente por la recalibración de umbrales y, en menor medida, por los \textit{prompts} anti-sesgo, como muestra el estudio de ablación de §10.2. Si en el piloto la tasa de falsos negativos supera el 2 \%, la acción inmediata sería bajar el umbral de 45 a un valor más conservador.

El análisis de sensibilidad muestra que la reducción de tiempo supera el 40 \% incluso en el escenario más pesimista (57 \%). Sin embargo, la reducción de tiempo depende de supuestos sobre el comportamiento del auditor que no han sido medidos, y la concordancia del \textit{Golden Set} (58.8 \%) revela una tasa de sobre-escalación que requiere calibración adicional. La validación operativa con un auditor real en producción queda pendiente; el MVP deja preparada la infraestructura para ello.

Los datos que sustentan esta afirmación son tres. Primero, el \textit{benchmark} de 505 evaluaciones demostró que Claude Sonnet 4 discrimina entre casos de bajo y alto riesgo con una distribución equilibrada ($\sigma$ = 20.1), lo que habilita la clasificación automática. Segundo, el \textit{pipeline} integrado (Fase 1 + Fase 2 con estrategia escalonada) procesó las 101 liquidaciones sin errores, produciendo 66 fast-track (65.3 \%), 9 revisiones rápidas (8.9 \%) y 26 auditorías completas (25.7 \%), con un costo total de API de USD 12.00 en 38.9 minutos. Tercero, el \textit{Golden Set} de 17 casos, diseñado como herramienta de regresión y calibración y no como conjunto de validación independiente, mostró una concordancia de 10/17 (58.8 \%) con el \textit{pipeline} real. Los 7 \textit{mismatches} son sobre-escalaciones (falsos positivos, no falsos negativos), lo que indica una dirección segura del error pero también un margen de calibración pendiente.

Lo que falta es la validación con el usuario final. La métrica de concordancia auditor-sistema y la tasa de falsos positivos en condiciones reales solo pueden obtenerse mediante un piloto controlado. El MVP deja preparada la infraestructura para esa medición (interfaz con registro de acciones, \textit{logging} de \textit{pipeline}), pero el dato no existe aún.

\section{Contribución}

\textbf{\textit{NexusCare}} aporta dos elementos concretos al estado del arte de la auditoría médica automatizada en el Perú. Lo más transferible es la demostración de que la arquitectura de dos fases (reglas + LLM) complementa al LLM único en una dimensión que el LLM solo no cubre: la verificación contractual determinística. El estudio de ablación (§10.2) mostró que Fase 1 no mejora la tasa de fast-track (el LLM con \textit{prompts} calibrados alcanza un 72.3 \% sin contexto contractual), pero sí mejora la calidad de la discriminación al escalar los 7 \textit{claims} con discrepancias contractuales que el LLM ignoraba. En dominios donde conviven reglas de negocio explícitas con juicios que requieren razonamiento sobre lenguaje natural, la combinación no es aditiva sino complementaria: las verificaciones contractuales son determinísticas y reproducibles, mientras que el LLM se concentra en lo que genuinamente aporta, el razonamiento clínico contextual. El segundo es el benchmark comparativo de cinco modelos sobre datos reales del sector salud peruano: la superioridad del generalista sobre el especialista, la importancia de sigma como métrica de discriminación y la inviabilidad de GPT-5.4 por sobre-penalización son hallazgos empíricos reproducibles. No se encontró literatura previa que compare modelos de lenguaje en esta tarea específica, ni en el Perú ni en la región.

El diseño del \textit{pipeline} con \textit{privacy by design}, trazabilidad y \textit{Golden Set} para regresión continua complementa la contribución, aunque responde más a buenas prácticas de ingeniería que a una innovación metodológica.

\section{Viabilidad}

La viabilidad técnica está demostrada: el \textit{pipeline} procesa 101 liquidaciones sin errores fatales, los 190 \textit{tests} pasan, el costo de API es marginal frente al valor de las liquidaciones, y la latencia es compatible con procesamiento por lotes.

La viabilidad económica es favorable con reservas. El costo de API medido (USD 12.00 por lote con estrategia escalonada, 33 \% de escalación a Opus) es insignificante. El ahorro operativo estimado (de 75.8 horas a 21–32.5 horas según el escenario de adopción, una reducción del 57–72 \%) justifica la inversión en desarrollo y mantenimiento del sistema. La reserva es la dependencia del \textit{pricing} de Anthropic: un aumento significativo en los precios de API alteraría la ecuación, aunque el sistema soporta la migración a otros modelos (Gemini, por ejemplo, a USD 1.78 por lote).

La viabilidad operativa es la más incierta. Requiere tres condiciones que no dependen del equipo técnico: (a) disposición de una EPS/IAFAS a participar en un piloto con datos de producción, (b) aceptación del auditor médico del flujo de trabajo propuesto, y (c) claridad regulatoria de SUSALUD sobre el uso de IA como herramienta de soporte en la adjudicación de liquidaciones. Sin estas tres condiciones, \textbf{\textit{NexusCare}} permanece como un proyecto académico técnicamente validado.
